% 使用 XeLaTeX 编译；版式参照泛函分析作业。
\documentclass{article}
\usepackage[UTF8]{ctex}
\usepackage{sectsty}
\usepackage{amssymb}
\usepackage{amsmath}
\usepackage{amsthm}
\usepackage{mathrsfs}
\usepackage{graphicx}
\title{拓扑学作业（第二周·题目）}
\author{张溢致\quad 致理数41\quad 2024011936}
\date{}
\newcommand{\R}{\mathbb R}
\newcommand{\problem}[2]{\par\medskip\noindent\textbf{#1.}\ （#2）\quad}

\begin{document}
\maketitle

\noindent\textbf{提交时间：}2026 年 10 月 1 日上午 8:00 前。

\problem{1}{Armstrong，第 2 章，第 27 题}
设 $(X,d)$ 是度量空间，$A\subseteq X$。证明
\[
  d(x,A)=0 \quad\Longleftrightarrow\quad x\in\overline A,
\]
其中 $d(x,A)=\inf\{d(x,a):a\in A\}$。

\problem{2}{Armstrong，第 2 章，第 28 题}
设 $A,B$ 是度量空间 $(X,d)$ 中两个不相交的闭子集。求两个不相交的开集 $U,V$，使得
\[
  A\subseteq U,\qquad B\subseteq V.
\]

\problem{3}{Armstrong，第 2 章，第 29 题}
在任意集合 $X$ 上定义
\[
  d(x,y)=
  \begin{cases}
    0,&x=y,\\
    1,&x\ne y.
  \end{cases}
\]
证明 $d$ 是一个度量，并说明它在 $X$ 上诱导什么拓扑。

\problem{4}{Armstrong，第 2 章，第 30 题}
证明度量空间的每个闭子集都是可数多个开集的交。

\problem{5}{Armstrong，第 2 章，第 31 题}
若 $A,B$ 是度量空间的子集，定义它们之间的距离为
\[
  d(A,B)=\inf\{d(x,y):x\in A, y\in B\}.
\]
求平面中两个不相交的闭子集，使它们之间的距离为零。并验证所求两个闭集的直径均为无穷大，其中
\[
  \operatorname{diam}(A)=\sup\{d(x,y):x,y\in A\}.
\]

\problem{6}{Armstrong，第 2 章，第 35 题}
给定连续映射
\[
  f:X\longrightarrow\R^{n+1}\setminus\{0\},
\]
求连续映射 $g:X\to S^n$，使得 $g$ 在 $f^{-1}(S^n)$ 上与 $f$ 相同。

\problem{7}{Munkres，第 7 章，§43，第 4 题}
证明度量空间 $(X,d)$ 完备，当且仅当：对每个由非空闭集组成的递减序列
\[
  A_1\supseteq A_2\supseteq A_3\supseteq\cdots,
\]
只要 $\operatorname{diam}(A_n)\to0$，就有
\[
  \bigcap_{n=1}^{\infty}A_n\ne\varnothing.
\]

\end{document}
